% Options for packages loaded elsewhere
\PassOptionsToPackage{unicode}{hyperref}
\PassOptionsToPackage{hyphens}{url}
\PassOptionsToPackage{dvipsnames,svgnames,x11names}{xcolor}
%
\documentclass[
  11pt,
]{article}
\usepackage{amsmath,amssymb}
\usepackage{iftex}
\ifPDFTeX
  \usepackage[T1]{fontenc}
  \usepackage[utf8]{inputenc}
  \usepackage{textcomp} % provide euro and other symbols
\else % if luatex or xetex
  \usepackage{unicode-math} % this also loads fontspec
  \defaultfontfeatures{Scale=MatchLowercase}
  \defaultfontfeatures[\rmfamily]{Ligatures=TeX,Scale=1}
\fi
\usepackage{lmodern}
\ifPDFTeX\else
  % xetex/luatex font selection
  \setmainfont[]{Liberation Serif}
\fi
% Use upquote if available, for straight quotes in verbatim environments
\IfFileExists{upquote.sty}{\usepackage{upquote}}{}
\IfFileExists{microtype.sty}{% use microtype if available
  \usepackage[]{microtype}
  \UseMicrotypeSet[protrusion]{basicmath} % disable protrusion for tt fonts
}{}
\makeatletter
\@ifundefined{KOMAClassName}{% if non-KOMA class
  \IfFileExists{parskip.sty}{%
    \usepackage{parskip}
  }{% else
    \setlength{\parindent}{0pt}
    \setlength{\parskip}{6pt plus 2pt minus 1pt}}
}{% if KOMA class
  \KOMAoptions{parskip=half}}
\makeatother
\usepackage{xcolor}
\usepackage[margin=25mm]{geometry}
\setlength{\emergencystretch}{3em} % prevent overfull lines
\providecommand{\tightlist}{%
  \setlength{\itemsep}{0pt}\setlength{\parskip}{0pt}}
\setcounter{secnumdepth}{5}
\ifLuaTeX
  \usepackage{selnolig}  % disable illegal ligatures
\fi
\usepackage{bookmark}
\IfFileExists{xurl.sty}{\usepackage{xurl}}{} % add URL line breaks if available
\urlstyle{same}
\hypersetup{
  pdftitle={Exact 1-Perfect-Code Classification of the Tribonacci Cubes Q\_n(111)},
  pdfauthor={Ryutaro Yonezu --- Independent Researcher},
  pdfkeywords={Tribonacci cube, generalized Fibonacci cube, 1-perfect
code, efficient domination, parity obstruction},
  colorlinks=true,
  linkcolor={black},
  filecolor={Maroon},
  citecolor={Blue},
  urlcolor={blue},
  pdfcreator={LaTeX via pandoc}}

\title{Exact 1-Perfect-Code Classification of the Tribonacci Cubes
\(Q_n(111)\)}
\author{Ryutaro Yonezu --- Independent Researcher}
\date{17 November 2026}

\begin{document}
\maketitle
\begin{abstract}
Let \(Q_n(111)\) be the graph induced by all binary strings of length
\(n\) avoiding three consecutive ones. We establish that this graph
admits a 1-perfect code if and only if \(n=1\) or \(n=3\). The
nonexistence argument for all remaining dimensions uses low-weight
sphere-partition constraints, endpoint rigidity of a forced singleton
codeword, and a parity obstruction on a family of four-element vertices.
The result addresses the fixed forbidden word 111 in the generalized
Fibonacci cube family. This is an unreviewed, AI-assisted preprint. Scholarly novelty is inconclusive and references are partial.
\end{abstract}

\textbf{Status:} Unreviewed preprint, version 1.0. Manuscript date: 17 November 2026.
Independent AI audit reported no fatal proof defects within its scope;
formal proof checking and human peer review have not been performed.
Author-reported Windows finite tests PASS. Novelty INCONCLUSIVE;
references PARTIAL. No priority claim. Copyright (c) 2026 Ryutaro
Yonezu. All rights reserved.

\section{Definitions}\label{definitions}

For n \textgreater= 1, let \(V_n\) be all subsets S of {[}n{]} that do
not contain any of the consecutive triples \{i,i+1,i+2\}, 1 \textless= i
\textless= n-2. Equivalently, these are all length-n binary words
avoiding 111. In \(Q_n(111)\), two vertices are adjacent when their
supports differ by exactly one element. Since \(V_n\) is downward
closed, deleting a 1-bit always gives a legal vertex.

A 1-perfect code is a subset C of \(V_n\) such that each vertex has
\textbf{exactly one} codeword in its closed neighborhood (itself or a
one-bit flip). Distinct codeword closed neighborhoods do not overlap.

\section{Position relative to earlier
work}\label{position-relative-to-earlier-work}

The graph \(Q_n(111)\) is the generalized Fibonacci cube denoted
\(\Gamma_n(111)=\Gamma_n(1^3)\) in the cited literature. It is distinct
from the ordinary Fibonacci cube \(\Gamma_n(11)\) and from the cyclic
Lucas cube.

Mollard (2018) studied perfect codes in generalized Fibonacci
cubes\(\Gamma_n(1^s)\), namely binary words avoiding a run of s ones.
His Conjecture 3.1 states that, for \(n\ge 3\) and \(s\ge 1\),
\textbf{if} a perfect code C exists in \(\Gamma_n(1^s)\), then
\(n=2^p-1\) for some integer p \textbf{and furthermore C is a perfect
code in the full n-cube Q\_n}. The present fixed-\(s=3\) classification,
\textbf{if the proof survives further independent verification},
establishes this conjecture for \(s=3\): the only \(n\ge 3\) with a
perfect code is \(n=3\), and the exhibited C=\{100,011\} is also perfect
in \(Q_3\). The main theorem also handles \(n=1\) outside Mollard's
\(n\ge 3\) hypothesis.

Mollard's existence corollary covers \(n=2^p-1\) for
\(s\ge 3\cdot 2^{p-2}\), giving \(s=3\) for \(p=2\) (\(n=3\)) but not for
\(p\ge 3\). His discussion gives a code for \(n=7\), \(s=5\); this does
not decide \(n=7\), \(s=3\). The 2018 paper consequently does not
subsume the fixed-\(s=3\) nonexistence statement here. This is a
\textbf{scope comparison, not a comprehensive novelty or priority
determination}.

References: Michel Mollard, \emph{The existence of perfect codes in a
family of generalized Fibonacci cubes}, Information Processing Letters
(2018), DOI: 10.1016/j.ipl.2018.07.010, arXiv:1801.04106; A. R. Ashrafi
et al., \emph{The (non-)existence of perfect codes in Fibonacci cubes},
Information Processing Letters (2016), DOI: 10.1016/j.ipl.2016.01.010.

\section{Main theorem}\label{main-theorem}

\textbf{For every integer n \textgreater= 1, \(Q_n(111)\) admits a
1-perfect code if and only if \(n=1\) or \(n=3\).} For \(n=1\), use
C=\{0\}. For \(n=3\), use C=\{100,011\}; their closed neighborhoods
partition exactly the seven legal words. We prove nonexistence
otherwise.

\subsection{\texorpdfstring{Lemma 1. For \(n\ge 4\), the empty set is
not a
codeword.}{Lemma 1. For n\textbackslash ge 4, the empty set is not a codeword.}}\label{lemma-1.-for-nge-4-the-empty-set-is-not-a-codeword.}

Suppose it were. Then all singleton vertices are dominated by the empty
set, so none of them is a codeword. Nor can a pair \{a,b\} be a
codeword: the singleton \{a\} would then have two dominating centers,
the empty word and \{a,b\}. Every pair of {[}n{]} must therefore be
dominated by a triple codeword containing it. Every triple in \(V_n\) is
legal only if it is not a consecutive triple; it nevertheless has all
three pairs as legal lower neighbors. The triple codewords form a
Steiner triple system on the n coordinates. If n is not congruent to 1
or 3 modulo 6, a Steiner triple system cannot exist, giving the required
contradiction immediately. It remains to consider n congruent to 1 or 3
modulo 6, when the system would have \(n(n-1)/6\) triple codewords.

Exactly n-2 of all binom(n,3) triples are forbidden, so
\(N_3=\binom{n}{3}-(n-2)\). Every legal triple not itself a codeword
must be dominated by a four-element codeword: there are no pair
codewords, and a center of equal weight cannot be adjacent to a
different triple. Each legal four-element codeword dominates exactly
four legal triples by deletion of one coordinate. Thus

\[4\mid N_3-\frac{n(n-1)}{6}=\frac{n(n-1)(n-3)}{6}-(n-2).\]

For \(n=1\) or 3 modulo 6, the first term on the right is even, while
n-2 is odd. The right side is odd, contradiction. Hence the empty word
is not a codeword. QED.

\subsection{\texorpdfstring{Lemma 2. For \(n\ge 4\), a code would have
odd n, one singleton center e\_j, and a matching M of pair codewords on
{[}n{]}\textbackslash\{j\}.}{Lemma 2. For n\textbackslash ge 4, a code would have odd n, one singleton center e\_j, and a matching M of pair codewords on {[}n{]}\textbackslash\{j\}.}}\label{lemma-2.-for-nge-4-a-code-would-have-odd-n-one-singleton-center-e_j-and-a-matching-m-of-pair-codewords-on-nj.}

By Lemma 1, the empty vertex has to be dominated by precisely one
singleton center e\_j. Every other singleton e\_a (a != j) cannot be a
center and must be dominated by a pair codeword \{a,b\}. (No pair
containing j can be a center because it is adjacent to e\_j.) Each pair
codeword dominates exactly two such singleton vertices, and each of the
n-1 remaining singleton vertices has exactly one dominating pair.
Therefore these pair codewords form a perfect matching M of
{[}n{]}\textbackslash\{j\}, and n-1 is even. Thus n is odd. QED.

\subsection{\texorpdfstring{Lemma 3 (endpoint rigidity). For \(n\ge 5\),
j is 1 or n.~If \(j=1\), \{2,3\} is an edge of
M.}{Lemma 3 (endpoint rigidity). For n\textbackslash ge 5, j is 1 or n.~If j=1, \{2,3\} is an edge of M.}}\label{lemma-3-endpoint-rigidity.-for-nge-5-j-is-1-or-n.-if-j1-23-is-an-edge-of-m.}

For every forbidden consecutive triple containing j, delete j and retain
the other two indices. These pairs form an edge set F\_j on
X={[}n{]}\textbackslash\{j\}. A direct inspection of consecutive triples
shows that \(F_j\) is a path of length 1 at an endpoint \(j=1\) or j=n,
length 2 at \(j=2\) or j=n-1, and length 3 at all remaining interior j.
In particular, \(F_j\) is a forest, and for an interior j it contains
two adjacent edges.

Fix a in X. Consider legal three-element vertices \{j,a,b\} with b in
X\textbackslash\{a\} that are not dominated by the unique pair codeword
of M incident to a. There are exactly

\[n-3-\deg_{F_j\setminus M}(a)\]

such vertices: of the n-2 choices of b, delete the forbidden choices
(edges of F\_j incident with a), then delete the matching partner unless
it was already forbidden. None of these triples can be a codeword (any
codeword containing j and two other elements shares a
closed-neighborhood vertex with e\_j). No other pair codeword can
dominate them. Hence they are dominated by four-element codewords
containing \{j,a\}. Every such four-element center dominates exactly
\emph{two} triples of this family, by deleting each of its other two
elements. Downward closure makes both triples legal; neither can be
dominated by the matching pair center as well, since that would violate
unique domination. Thus both are among the triples being counted.
Therefore \(n-3-\deg_{F_j\setminus M}(a)\) is even.

As n is odd, \(\deg_{F_j\setminus M}(a)\) is even for every a. A
nonempty forest contains an odd-degree vertex. Hence \(F_j\setminus M\)
must be empty and \(F_j\) is contained in M. If j were interior, \(F_j\)
would contain adjacent edges, impossible in a matching. Therefore j is
an endpoint. In the \(j=1\) case, \(F_1=\{\{2,3\}\}\), so \{2,3\} is in
M. QED.

\subsection{\texorpdfstring{Lemma 4 (odd fan obstruction). For every odd
\(n\ge 5\), the endpoint configuration of Lemma 3 is
impossible.}{Lemma 4 (odd fan obstruction). For every odd n\textbackslash ge 5, the endpoint configuration of Lemma 3 is impossible.}}\label{lemma-4-odd-fan-obstruction.-for-every-odd-nge-5-the-endpoint-configuration-of-lemma-3-is-impossible.}

By reversal, suppose \(j=1\). Lemma 3 forces the singleton center \{1\}
and pair center \{2,3\}. Let P=\{2,n\} and S=\{1,2,n\}. Both vertices
are legal; P is not a pair codeword because 2 is already matched to 3.
As no singleton center dominates P, its unique dominating center is a
triple R=P union \{c\} for some c in L=\{4,5,\ldots,n-1\}. The index c
cannot be 3, since \{2,3\} is already a codeword and would be adjacent
to \{2,3,n\}.

The legal triple S contains the singleton-center coordinate 1, so S
cannot be a codeword: for instance the vertex \{1,2\} would lie in the
closed neighborhoods of both S and the singleton codeword \{1\}. More
generally, no triple containing 1 can be a codeword, since deleting
either of its other coordinates produces a legal pair adjacent to both
that triple and \{1\}. It is not dominated by a pair center: none of its
pairs is a codeword (the two pairs containing 1 overlap the singleton
center, and P is not in M). Hence its unique dominating center is a
four-element word D=S union \{d\}, with d in L. The index 3 is not
allowed, since S union \{3\} contains 111 at positions 1,2,3. The
centers R and D must be distinct and have disjoint closed neighborhoods.
In particular c != d, because if c=d, R and D are adjacent. \textbf{For
\(n=5\), the set L=\{4\} has only one element, so distinct c,d in L
cannot exist: this already gives a contradiction.} Hence below it
suffices to consider odd \(n\ge 7\).

Define the family of legal four-element vertices

\[W_x=S\cup\{x\},\quad x\in L=\{4,\ldots,n-1\}.\]

There are \(|L|=n-4\) of these. Exactly one is dominated by a codeword
of weight 3: any triple codeword dominating W\_x must be a three-element
subset, and as no triple codeword contains 1 (by the singleton \{1\}),
its support must be \{2,n,x\}. Since P has exactly one dominating triple
R, this occurs exactly when x=c.~Exactly one W\_x is a codeword of
weight 4: any such W\_x would also dominate S, which has unique
dominating codeword D, so W\_x=D, occurring exactly when x=d.~Weight-2
or smaller centers cannot dominate weight-4 words, and a distinct
weight-4 center cannot be adjacent to a weight-4 word in a binary cube
(equal-weight words differ in an even number of bits).

Every remaining W\_x must thus be dominated by a five-element center.
Such a center is necessarily S union \{x,z\}, with z in L distinct from
x: all legal centers contain S, the bit z=3 is forbidden, and the other
positions in S are already present. A five-element center of this form
dominates exactly two vertices of the family, namely W\_x and W\_z. As
the closed neighborhoods of codewords partition the vertices, no such
five-element center can also dominate either of the two already covered
words W\_c and W\_d.~Consequently, for odd \(n\ge 7\), the \(n-6\) other
words must be partitioned into pairs. Thus n-6 is even, implying n even,
contrary to Lemma 2. Combined with the preceding \(n=5\) contradiction,
this completes the lemma. QED.

\subsection{Completion}\label{completion}

For even \(n\ge 4\), Lemma 2 gives a contradiction. For odd \(n\ge 5\),
Lemmas 2--4 give a contradiction. At \(n=2\), the graph is \(Q_2\), a
4-cycle whose closed balls each have size 3, so a perfect partition is
impossible. Finally, the explicit codes at \(n=1\) and \(n=3\) give
existence. Therefore the theorem holds for every \(n\ge 1\). QED.

\section{Verification and research
status}\label{verification-and-research-status}

The mathematical statements, Lemmas 1--4 and their proofs in the Main
theorem through Completion section above are reproduced verbatim from
frozen No.71 v0.7. The original author-provided archive and its SHA256
identify the mathematical source. A Python standard-library finite
sanity checker is supplied separately; it checks small cases and local
identities, not the infinite-dimensional theorem. The author reported a
successful Windows PowerShell run of the v0.5 verifier. An external AI
auditor reported an independent SAT check through \(n=17\), but neither
that SAT implementation nor its certificates are part of the present
public reproduction package. Human peer review and proof-assistant
formal verification have not been performed. Scholarly novelty remains
INCONCLUSIVE and literature coverage PARTIAL. No priority or world-first
claim is made.

\section{References}\label{references}

\begin{enumerate}
\def\labelenumi{\arabic{enumi}.}
\tightlist
\item
  M. Mollard, \emph{The existence of perfect codes in a family of
  generalized Fibonacci cubes}. Information Processing Letters (2018).
  DOI: https://doi.org/10.1016/j.ipl.2018.07.010. arXiv:1801.04106.
\item
  A. R. Ashrafi, J. Azarija, A. Babai, K. Fathalikhani and S. Klavžar,
  \emph{The (non-)existence of perfect codes in Fibonacci cubes}.
  Information Processing Letters \textbf{116} (2016), 387--390. DOI:
  https://doi.org/10.1016/j.ipl.2016.01.010.
\item
  Y. Rho and A. Vesel, \emph{Linear recognition of generalized Fibonacci
  cubes \(Q_h(111)\)}. Discrete Mathematics \& Theoretical Computer
  Science \textbf{17} (2016), 17:3. DOI:
  https://doi.org/10.46298/dmtcs.2165.
\end{enumerate}

\end{document}
